\section*{Введение}
\addcontentsline{toc}{section}{Введение}

\textbf{Актуальность темы.} Рынок браузерных игр демонстрирует устойчивый рост:
по данным аналитиков, аудитория HTML5-игр превысила 3 миллиарда пользователей.
Независимые разработчики нуждаются в доступных инструментах публикации, однако
существующие платформы либо ориентированы на коммерческий сегмент со значительной
комиссией, либо не поддерживают современные форматы сборок, в частности Unity WebGL.
Разработка открытой веб-платформы, обеспечивающей полный цикл от загрузки игры
до её запуска в браузере, является актуальной практической задачей.

\textbf{Степень разработанности проблемы.} Существующие платформы --- itch.io,
GameJolt, Яндекс Игры --- решают задачу частично: itch.io не предоставляет
структурированной системы ролей и агрегации статистики, Яндекс Игры требуют
интеграции с проприетарным SDK и рекламной системой. Открытых полнофункциональных
решений на стеке Flask + PostgreSQL с поддержкой Unity WebGL в открытом доступе
не представлено.

\textbf{Объект исследования} --- процессы публикации, распространения и запуска
браузерных игр в рамках веб-платформы.

\textbf{Предмет исследования} --- программные средства и архитектурные решения
для реализации платформы с поддержкой загрузки игровых сборок, ролевой модели
и агрегации пользовательских данных.

\textbf{Целью данной работы} является проектирование и реализация веб-платформы
WhyGame, обеспечивающей публикацию, хранение и браузерный запуск игр формата
Unity WebGL с поддержкой системы комментариев, рейтингов и каталога с фильтрацией.